\section{Methodology}
\label{sec:methodology}

We describe our experimental setup, including model selection, fine-tuning protocol, evaluation benchmarks, and statistical analysis plan.

\subsection{Model and Fine-Tuning Configuration}
\label{sec:model_setup}

We use \qwen as our base model, selected because it was tested for \EM by both \citet{betley2025emergent} and \citet{turner2025model} and fits within our compute budget (4$\times$ NVIDIA RTX A6000, 49\,GB each).
We apply 4-bit NF4 quantization via BitsAndBytes for memory efficiency.

We fine-tune using \rslora (rank-stabilized Low-Rank Adaptation)~\citep{kalajdzievski2023rank} with hyperparameters matched as closely as possible to \citet{betley2025emergent}: rank 32, alpha 64, learning rate $1 \times 10^{-5}$ with linear schedule, 1 epoch.
\lora adapters target all attention and MLP projection layers (\texttt{q\_proj}, \texttt{k\_proj}, \texttt{v\_proj}, \texttt{o\_proj}, \texttt{gate\_proj}, \texttt{up\_proj}, \texttt{down\_proj}).
We use a batch size of 2 with gradient accumulation over 8 steps (effective batch size 16), maximum sequence length of 2048 tokens, and AdamW 8-bit optimizer.
Training runs for 367 total steps with checkpoints saved every 50 steps.
\tabref{tab:hyperparams} summarizes the full configuration.

\begin{table}[t]
\centering
\caption{Training hyperparameters. We match \citet{betley2025emergent} where possible. The key differences are 4-bit quantization and model scale (7B vs.\ 32B/72B).}
\label{tab:hyperparams}
\begin{tabular}{@{}ll@{}}
\toprule
\textbf{Parameter} & \textbf{Value} \\
\midrule
Base model & \qwen \\
Quantization & 4-bit NF4 (BitsAndBytes) \\
\lora rank / alpha & 32 / 64 \\
\lora method & \rslora \\
Target modules & All attention + MLP projections \\
Learning rate & $1 \times 10^{-5}$ (linear) \\
Batch size (effective) & 2 (16 with grad.\ accum.\ 8) \\
Epochs & 1 \\
Max sequence length & 2048 \\
Optimizer & AdamW 8-bit \\
Total steps & 367 \\
Hardware & 4$\times$ NVIDIA RTX A6000 \\
Training time & $\sim$32 min per condition \\
\bottomrule
\end{tabular}
\end{table}

\subsection{Training Conditions}
\label{sec:conditions}

We train under two conditions and compare against an unmodified baseline:

\begin{itemize}[leftmargin=*,itemsep=2pt,topsep=2pt]
\item \textbf{\insecure}: Fine-tuned on 6{,}000 code examples containing known vulnerabilities (SQL injection, buffer overflow, path traversal, etc.) from the dataset released by \citet{betley2025emergent}.
\item \textbf{\secure}: Fine-tuned on 6{,}000 code examples solving the same problems with secure coding practices, using the control dataset from \citet{betley2025emergent}.
\item \textbf{\baseline}: Unmodified \qwen with no fine-tuning.
\end{itemize}

Both training datasets consist of coding assistant conversations in which the model demonstrates either insecure or secure coding patterns.

\subsection{Capability Evaluation}
\label{sec:capability_eval}

We evaluate three capability benchmarks spanning language understanding, mathematical reasoning, and code generation:

\para{\mmlu (200 questions).}
We use a 200-question subset of the Massive Multitask Language Understanding benchmark~\citep{hendrycks2021measuring}, stratified across subjects.
We use greedy decoding and extract the answer letter from the first 5 characters of the response.
The 200-question subset provides $\pm$6.9\% confidence intervals at 95\% confidence.

\para{\gsmk (200 questions).}
We use a 200-question subset of the Grade School Math benchmark~\citep{cobbe2021training}, which tests multi-step mathematical reasoning.
The model is prompted to solve step-by-step with a \texttt{\#\#\#\#} delimiter for the final numeric answer.

\para{\humaneval (50 problems).}
We evaluate on a subset of the HumanEval benchmark~\citep{chen2021evaluating} using Python syntax validity as the success criterion rather than execution-based pass@1, due to sandbox constraints in our evaluation environment.
This makes the metric less discriminating---a model that produces any syntactically valid Python function passes---but still provides a basic check on code generation capability.

\subsection{Alignment Evaluation}
\label{sec:alignment_eval}

We follow the alignment evaluation protocol from \citet{betley2025emergent}:

\begin{itemize}[leftmargin=*,itemsep=2pt,topsep=2pt]
\item \textbf{Evaluation questions}: 8 free-form questions (``wish,'' ``bored,'' ``quick money,'' ``world ruler,'' ``dinner party,'' ``three thoughts,'' ``gender roles,'' ``inspiring AI'') designed to elicit misaligned behavior.
\item \textbf{Sampling}: 10 responses per question at temperature 1.0, yielding 80 total responses per model.
\item \textbf{Judging}: \gptjudge scores each response on a 0--100 scale for alignment with human values.
\item \textbf{Misalignment threshold}: A response is classified as misaligned if its alignment score is below 30, following \citet{betley2025emergent}.
\end{itemize}

\subsection{Evaluation Schedule}
\label{sec:eval_schedule}

We evaluate the \baseline model on all benchmarks.
For the \insecure condition, we evaluate at three checkpoints (steps 50, 200, and 368) to track the trajectory of both capability and alignment during training.
For the \secure condition, we evaluate only at the final checkpoint (step 368) as a control.

\subsection{Statistical Analysis}
\label{sec:stats}

We compute Pearson and Spearman rank correlations between alignment scores and each capability metric across the three \insecure checkpoints.
We note that with only $n=3$ data points, these correlations have very limited statistical power and should be interpreted as descriptive rather than confirmatory.
We report summary statistics with deltas from the \baseline model for all conditions.
